\section{An upper bound and the dependence on ambient dimension}
\label{sec:upper-bound}
An upper bound identifies the natural comparison scale in the sparse
regime.

\begin{proposition}[An actual-posterior upper bound]\label{prop:upper-bound}
Let $\lambda_k$ be the normalized Walsh eigenvalue defined above.
There is an absolute constant $C$ such that, for every
$0<\varepsilon<1/3$, $0<\delta<1$, and
\[
 n\ge \frac{C}{\varepsilon^2}
       \left(\lambda_k^{-1}+\log\frac2\delta\right),
\]
the actual posterior at every fixed $0<\gamma<1$ satisfies
\[
 \mathbb P_S\left\{\limsup_{N\to\infty}F_N(S)
                      \le3\varepsilon\right\}\ge1-\delta.
\]
On the same sample event, every compatible minimum-energy sign function
has test error at most $\varepsilon$.
\end{proposition}
\begin{proof}
The globally separated comparator $z=af$ is feasible for every sample
and has $\|af\|_K^2=a^2/\lambda_k$. Thus every minimizing field lies
in the RKHS ball of radius $R=a/\sqrt{\lambda_k}$ and satisfies
$f(X_i)z(X_i)\ge a$.
Use the ramp loss
$\ell(t)=\min\{1,\max\{0,1-t/a\}\}$.
Its empirical loss vanishes for these fields, and
$\mathbf1_{\{f(x)z(x)\le0\}}\le\ell(f(x)z(x))$.

For completeness, the standard uniform deviation estimate here is
\[
 \sup_{\|z\|_K\le R}
 \left[\mathbb E\ell(fz)-\frac1n\sum_i\ell(f(X_i)z(X_i))\right]
 \le \frac{4R}{a\sqrt n}
       +\sqrt{\frac{\log(1/\delta)}{2n}}
\]
with probability at least $1-\delta$.
Indeed, symmetrization and contraction for the $a^{-1}$-Lipschitz
function $\ell-\ell(0)$ bound the expected supremum by $4/a$ times
the RKHS Rademacher complexity. The reproducing identity and
$K(x,x)=1$ give
\[
 \frac1n\mathbb E_\sigma
 \sup_{\|z\|_K\le R}\left|\sum_i\sigma_i f(X_i)z(X_i)\right|
 =\frac Rn\mathbb E_\sigma
       \left\|\sum_i\sigma_i f(X_i)K_{X_i}\right\|_K
 \le\frac R{\sqrt n}.
\]
Changing one sample changes the supremum deviation by at most $1/n$;
the bounded-differences inequality supplies the displayed probability
bound. The assumed sample size therefore makes every minimizing
sign's error at most $\varepsilon$.

By Proposition~\ref{prop:rate}, the actual posterior mass outside these
minimizing signs tends to zero. Hence its average test error
$\mathbb E_x[1-q_{N,S}(x)]$ has upper limit at most $\varepsilon$.
Since $F_N(S)\le3\mathbb E_x[1-q_{N,S}(x)]$, the claimed confidence
bound follows, including the strict greater-than-$2/3$ convention.
\end{proof}

For growing $k=o(d)$, put $j=k$ if $k$ is odd and $j=k+1$ otherwise.
Lemma~\ref{lem:sharp} gives
\[
 \lambda_k^{-1}\asymp j^{1+\beta}\binom{d+1}{j},
 \qquad
 \log\lambda_k^{-1}=\Theta\!\left(k\log\frac{ed}{k}\right).
\]
In particular, when $k=\lfloor d/\log d\rfloor$, fixed
$\varepsilon,\delta>0$ require only
$\exp\{O(k\log\log d)\}=d^{o(k)}$ samples in this upper bound.
Thus a literal $d^{\Omega(k)}$ lower bound, with a fixed positive
exponent uniformly over all growing sparse $k=o(d)$, is false for the
actual posterior. The upper bound has the natural exponential scale
$\exp\{\Theta(k\log(ed/k))\}$. A lower bound on this scale remains
open. For $k\le d^{1-\eta}$ with fixed $\eta>0$, such a lower bound
would entail $d^{\Omega(k)}$.
